%%%=====================================================================%%%
%%%                                                                      %%%
%%%   MATHÉMATIQUES POUR L'INTELLIGENCE ARTIFICIELLE                     %%%
%%%   Le tout-en-un des fondamentaux, accessible à tous                  %%%
%%%   Doctor Mokira — projet ML4E                                        %%%
%%%                                                                      %%%
%%%   Fichier maître. Compilation :                                      %%%
%%%       latexmk    (recommandé, tout automatique)                      %%%
%%%     ou :  make                                                       %%%
%%%     ou à la main :                                                   %%%
%%%       pdflatex -shell-escape main                                    %%%
%%%       bibtex   main                                                  %%%
%%%       pdflatex -shell-escape main                                    %%%
%%%       pdflatex -shell-escape main                                    %%%
%%%                                                                      %%%
%%%   IMPORTANT : l'option -shell-escape est indispensable (minted).     %%%
%%%=====================================================================%%%
\documentclass{ml4ebook}     % ajoutez l'option [brouillon] pour un filigrane

\begin{document}

%----------------------------------------------------------------------
% Pages liminaires
%----------------------------------------------------------------------
\frontmatter
%%%=====================================================================%%%
%%%  frontmatter/cover.tex — Page de titre                              %%%
%%%                                                                      %%%
%%%  Le titre, le sous-titre et l'auteur viennent des macros \mbTitre,   %%%
%%%  \mbSousTitre et \mbAuteur définies en tête de ml4ebook.cls.         %%%
%%%                                                                      %%%
%%%  PARTI PRIS : ce livre est un livre de MATHÉMATIQUES, dont l'IA est  %%%
%%%  le champ d'application — pas l'inverse. La couverture montre donc   %%%
%%%  un panorama des cinq domaines traités, chacun avec sa figure ET sa  %%%
%%%  formule, l'apprentissage n'étant que le dernier des cinq :          %%%
%%%                                                                      %%%
%%%     algèbre linéaire · analyse · probabilités                        %%%
%%%              optimisation · apprentissage                            %%%
%%%                                                                      %%%
%%%  Composition en registres horizontaux :                              %%%
%%%     1. la rangée haute des trois domaines fondamentaux               %%%
%%%     2. le bandeau bleu qui porte le titre                            %%%
%%%     3. le filet ambre, le sous-titre, la promesse                    %%%
%%%     4. la rangée basse : optimisation et apprentissage               %%%
%%%     5. l'auteur                                                      %%%
%%%                                                                      %%%
%%%  Les teintes sont NETTEMENT plus soutenues que dans le corps du      %%%
%%%  livre : une couverture est vue de loin, souvent imprimée, parfois   %%%
%%%  en noir et blanc. Un trait à 30 % d'opacité y disparaît.            %%%
%%%=====================================================================%%%
\begin{titlepage}
\thispagestyle{empty}

% Styles locaux : ils ne servent qu'ici, mais évitent de répéter dix fois
% les mêmes réglages de police pour les formules et les intitulés.
\tikzset{
  cov formule/.style={font=\small, text=mbInk},
  cov domaine/.style={font=\sffamily\scriptsize\bfseries, text=mbSlate},
  cov courbe/.style={draw=mbBlue, line width=1.1pt},
  cov axe/.style={draw=mbSlate!75, line width=0.6pt,
                  -{Stealth[round,length=1.7mm,width=1.3mm]}},
}

\begin{tikzpicture}[remember picture,overlay,mbfig]

  %--------------------------------------------------------------------
  % Fond
  %--------------------------------------------------------------------
  \fill[mbPaper] (current page.south west) rectangle (current page.north east);

  %====================================================================
  % 1. RANGÉE HAUTE — les trois piliers mathématiques
  %====================================================================
  \coordinate (rhaut) at ($(current page.north)+(0cm,-4.6cm)$);

  %--------------------------------------------------------------------
  % 1a. ALGÈBRE LINÉAIRE — une matrice, avec ses coefficients et ses
  %     crochets. Les cases colorées la rendent lisible de loin, les
  %     nombres disent que c'est bien un objet mathématique.
  %--------------------------------------------------------------------
  \begin{scope}[shift={($(rhaut)+(-4.5cm,0cm)$)}]
    \foreach \i in {0,...,3}{
      \foreach \j in {0,...,3}{
        \pgfmathtruncatemacro{\val}{mod(3*\i+5*\j+7,9)-4}
        \ifnum\i=\j
          \fill[mbAmber!45] (-1.16+\j*0.58, 1.16-\i*0.58) rectangle ++(0.58,-0.58);
        \else
          \fill[mbBlue!\the\numexpr8+4*\i+5*\j\relax]
                (-1.16+\j*0.58, 1.16-\i*0.58) rectangle ++(0.58,-0.58);
        \fi
        \draw[mbBlue!55, line width=0.4pt]
              (-1.16+\j*0.58, 1.16-\i*0.58) rectangle ++(0.58,-0.58);
        \node[font=\scriptsize, text=mbInk]
              at (-0.87+\j*0.58, 0.87-\i*0.58) {$\val$};
      }
    }
    % Crochets de matrice
    \draw[mbInk, line width=1pt]
      (-1.32,1.16) -- (-1.32,-1.16) (-1.32,1.16) -- (-1.10,1.16)
      (-1.32,-1.16) -- (-1.10,-1.16);
    \draw[mbInk, line width=1pt]
      (1.32,1.16) -- (1.32,-1.16) (1.32,1.16) -- (1.10,1.16)
      (1.32,-1.16) -- (1.10,-1.16);
    \node[cov formule] at (0,-1.78) {$A\vect{x} = \vect{b}$};
    \node[cov domaine] at (0,-2.28) {ALGÈBRE LINÉAIRE};
  \end{scope}

  %--------------------------------------------------------------------
  % 1b. ANALYSE — une courbe et sa tangente au sommet. La tangente y est
  %     horizontale : c'est exactement la condition f'(x) = 0, celle qui
  %     fonde toute l'optimisation. La figure annonce la suivante.
  %--------------------------------------------------------------------
  \begin{scope}[shift={($(rhaut)+(0cm,0cm)$)}]
    \draw[cov axe] (-1.98,-0.78) -- (2.05,-0.78);
    \draw[cov axe] (-1.8125,-0.95) -- (-1.8125,1.02);
    \draw[cov courbe, smooth, samples=80, domain=0:2.9]
      plot ({\x*1.25-1.8125}, {(0.75+0.5*sin(deg(1.9*\x)))*1.25-0.78});
    % La tangente au maximum : elle est HORIZONTALE, c'est-à-dire f'(x)=0
    \draw[draw=mbAmber!85!black, line width=1.2pt] (-1.5,0.7825) -- (-0.05,0.7825);
    \draw[mbSlate!55, line width=0.5pt, dashed] (-0.78,0.7825) -- (-0.78,-0.78);
    \fill[mbInk] (-0.78,0.7825) circle (2pt);
    \node[cov formule] at (0,-1.78) {$f'(x) = 0$};
    \node[cov domaine] at (0,-2.28) {ANALYSE};
  \end{scope}

  %--------------------------------------------------------------------
  % 1c. PROBABILITÉS — la courbe en cloche, avec l'aire sous la courbe :
  %     l'image même d'une probabilité.
  %--------------------------------------------------------------------
  \begin{scope}[shift={($(rhaut)+(4.5cm,0cm)$)}]
    % L'aire grisée court sur un écart-type de part et d'autre de la
    % moyenne : c'est une probabilité, montrée comme une surface.
    \fill[mbBlue!35]
      (-0.7375,-0.78) --
      plot[smooth, samples=50, domain=0.91:2.09]
        ({\x*1.25-1.875}, {1.25*exp(-(\x-1.5)^2/0.7)*1.25-0.78}) --
      (0.7375,-0.78) -- cycle;
    \draw[cov axe] (-2.05,-0.78) -- (2.1,-0.78);
    \draw[cov courbe, smooth, samples=80, domain=0:3]
      plot ({\x*1.25-1.875}, {1.25*exp(-(\x-1.5)^2/0.7)*1.25-0.78});
    \draw[mbAmber!85!black, line width=0.9pt, dashed] (0,-0.78) -- (0,0.78);
    \node[font=\scriptsize, text=mbAmber!80!black, below] at (0,-0.81) {$\mu$};
    \node[cov formule] at (0,-1.78) {$\Normale(\mu,\sigma^{2})$};
    \node[cov domaine] at (0,-2.28) {PROBABILITÉS};
  \end{scope}

  %====================================================================
  % 2. Bandeau du titre
  %    ATTENTION : les deux coins doivent avoir des ordonnées DIFFÉRENTES.
  %    La version d'origine allait de `current page.west` à
  %    `current page.east|-current page.center` — or `west` est déjà à la
  %    hauteur du centre : le rectangle était plat, donc invisible, et le
  %    titre blanc se retrouvait posé sur le fond crème, illisible.
  %====================================================================
  \fill[mbBlue] (current page.west|-current page.center)
        rectangle ([yshift=3.7cm]current page.east|-current page.center);
  \fill[mbAmber] (current page.west|-current page.center)
        rectangle ([yshift=-0.28cm]current page.east|-current page.center);

  \node[anchor=west,text width=11.5cm,align=left]
        at ([yshift=1.85cm,xshift=1.9cm]current page.west) {%
    \sffamily\bfseries\color{white}\fontsize{34}{39}\selectfont
    Machine Learning\\ For Engineer};

  %====================================================================
  % 3. Sous-titre et promesse, sous le filet — donc sur le fond crème
  %====================================================================
  \node[anchor=north west,text width=12.4cm,align=left]
        at ([yshift=-1.05cm,xshift=1.9cm]current page.west) {%
    \sffamily\color{mbInk}\large \mbSousTitre};
  \node[anchor=north west,text width=12.4cm,align=left]
        at ([yshift=-2.15cm,xshift=1.9cm]current page.west) {%
    \sffamily\bfseries\color{mbAmber!80!black}\normalsize
    Aucun prérequis : on part vraiment de zéro.};

  %====================================================================
  % 4. RANGÉE BASSE — ce qu'on construit avec les trois piliers
  %====================================================================
  \coordinate (rbas) at ($(current page.south)+(0cm,6.5cm)$);

  %--------------------------------------------------------------------
  % 4a. OPTIMISATION — lignes de niveau d'une fonction de coût et la
  %     descente qui rejoint le minimum. Le zigzag n'est pas décoratif :
  %     c'est le comportement réel d'une descente dans une vallée
  %     allongée, l'amplitude décroissant à chaque pas.
  %--------------------------------------------------------------------
  \begin{scope}[shift={($(rbas)+(-3.4cm,0cm)$)}]
    \begin{scope}[rotate=20]
      \foreach \r/\t in {1.15/34, 0.92/44, 0.69/54, 0.46/64, 0.23/74}{
        \draw[mbBlue!\t, line width=0.8pt]
              (0,0) ellipse ({\r*1.9cm} and {\r*1.08cm});
      }
      \fill[mbBlue] (0,0) circle (2pt);
      \draw[-{Stealth[round,length=2.2mm,width=1.7mm]},
            draw=mbAmber!85!black, line width=1.1pt]
        (-1.95,0.55) -- (-1.03,-0.46) -- (-0.55,0.32)
        -- (-0.27,-0.18) -- (-0.11,0.10) -- (-0.02,-0.01);
      \foreach \p in {(-1.95,0.55), (-1.03,-0.46), (-0.55,0.32),
                      (-0.27,-0.18), (-0.11,0.10)}{
        \fill[mbAmber!85!black] \p circle (1.8pt);
      }
    \end{scope}
    \node[cov formule] at (0,-1.85)
         {$\params \leftarrow \params - \pas\,\grad\cout(\params)$};
    \node[cov domaine] at (0,-2.35) {OPTIMISATION};
  \end{scope}

  %--------------------------------------------------------------------
  % 4b. APPRENTISSAGE — le réseau de neurones. Il vient en dernier, et
  %     il n'est qu'une application de tout ce qui précède : une
  %     multiplication par une matrice, puis une fonction, le tout réglé
  %     par une descente de gradient.
  %--------------------------------------------------------------------
  % `minimum size` doit être posé APRÈS le style `neurone`, dans la liste
  % d'options du nœud : passé par `every node`, il serait écrasé par le
  % 7 mm du style — plus large que l'espacement, et les neurones se
  % chevaucheraient.
  \begin{scope}[shift={($(rbas)+(3.4cm,0cm)$)}]
    \foreach \i in {0,1,2}
      \node[neurone entree, draw=mbAzure, fill=mbAzure!25, line width=0.8pt,
            minimum size=4.4mm]
            (re\i) at (-1.5, {0.62*(1-\i)}) {};
    \foreach \i in {0,1,2,3}
      \node[neurone cache, draw=mbBlue, fill=mbBlue!20, line width=0.8pt,
            minimum size=4.4mm]
            (rc\i) at (0, {0.62*(1.5-\i)}) {};
    \foreach \i in {0,1}
      \node[neurone sortie, draw=mbAmber!85!black, fill=mbAmber!45,
            line width=0.8pt, minimum size=4.4mm]
            (rs\i) at (1.5, {0.62*(0.5-\i)}) {};
    \foreach \a in {0,1,2} \foreach \b in {0,1,2,3}
      \draw[mbSlate!45, line width=0.45pt] (re\a) -- (rc\b);
    \foreach \a in {0,1,2,3} \foreach \b in {0,1}
      \draw[mbSlate!45, line width=0.45pt] (rc\a) -- (rs\b);
    \node[cov formule] at (0,-1.85)
         {$\vect{y} = \sigma(W\vect{x} + \vect{b})$};
    \node[cov domaine] at (0,-2.35) {APPRENTISSAGE};
  \end{scope}

  %====================================================================
  % 5. Auteur
  %====================================================================
  \node[anchor=south west,text width=12.4cm,align=left]
        at ([yshift=1.5cm,xshift=1.9cm]current page.south west) {%
    \sffamily
    {\bfseries\color{mbInk}\large\mbAuteur}\\[6pt]
    {\color{mbBlue}\footnotesize\faIcon{calendar-alt}\ \the\year}};
\end{tikzpicture}
\end{titlepage}

% Verso de la page de titre : colophon
\thispagestyle{empty}
\null\vfill
{\footnotesize\sffamily\color{mbSlate}
\noindent Édition \the\year. Document composé avec \LaTeX{} (pdf\TeX), polices
Palatino (texte et mathématiques), TeX~Gyre~Heros (titres) et DejaVu~Sans~Mono
(code).
Figures et schémas réalisés avec \textsf{TikZ} et \textsf{pgfplots}, coloration
syntaxique assurée par \textsf{minted}/Pygments, encadrés par
\textsf{tcolorbox}, bibliographie par \textsf{biblatex}.\par
\medskip
\noindent Ce livre rassemble en un seul volume les mathématiques dont
l'intelligence artificielle et la science des données ont réellement besoin :
algèbre linéaire, analyse, probabilités, statistique et optimisation. Il est
écrit pour être lu sans aucun prérequis : on part de zéro, on explique avec
des mots avant d'écrire des symboles, et on n'avance jamais sans avoir
d'abord dit pourquoi.\par}
\clearpage

\chapter*{Avant-propos}
\addcontentsline{toc}{chapter}{Avant-propos}
\markboth{Avant-propos}{}

\begin{quote}
\itshape « Un livre de mathématiques ne devrait effrayer personne. »
\end{quote}

L'intelligence artificielle repose entièrement sur des mathématiques. Le
problème, c'est que les livres qui les enseignent supposent presque tous
que vous les connaissez déjà : ils vont vite, ils écrivent en symboles
avant d'avoir expliqué en mots, et ils supposent une aisance en
programmation que peu de lecteurs ont au départ. Résultat : beaucoup de
gens curieux, capables, motivés, s'arrêtent à la première page.

\textbf{Ici, on ne suppose presque rien.} Chaque notion est d'abord
racontée avec des mots simples, puis illustrée par une image ou un
dessin, puis seulement écrite en langage mathématique — et enfin mise en
pratique avec quelques lignes de code que vous pouvez exécuter vous-même.

\begin{enclair}
Notre promesse : si vous savez additionner, multiplier et lire un tableau
de nombres, vous pouvez comprendre les mathématiques de l'intelligence
artificielle. Le reste, on le construit ensemble, dans l'ordre, sans
sauter d'étape.
\end{enclair}

\section*{Pour qui est ce livre}

Pour toute personne qui veut appliquer les mathématiques ailleurs —
informatique, IA, traitement du signal, analyse de données — et qui part
d'un niveau nul, ancien ou fragile. Aucun diplôme n'est demandé. Le texte
est écrit pour rester lisible par un élève de début de collège lorsque le
sujet le permet, et pour ne jamais perdre le lecteur lorsque le sujet
devient réellement difficile.

C'est aussi un \emph{tout-en-un} : plutôt que d'aller chercher l'algèbre
linéaire dans un livre, les probabilités dans un deuxième et
l'optimisation dans un troisième — chacun avec ses notations et ses
habitudes —, vous trouverez ici l'ensemble, avec des notations
cohérentes d'un bout à l'autre.

\section*{Comment lire ce livre}

Le texte utilise quelques encadrés récurrents, faciles à reconnaître à
leur pictogramme :

\begin{itemize}
  \item \textbf{\faIcon{comments}~En clair} — l'idée expliquée avec des
        mots de tous les jours.
  \item \textbf{\faIcon{brain}~L'intuition} — l'image mentale à garder en
        tête.
  \item \textbf{\faIcon{robot}~À quoi ça sert en IA ?} — où cette notion
        apparaît vraiment dans un algorithme d'apprentissage.
  \item \textbf{\faIcon{chart-line}~Sur des données réelles} — la même
        idée, appliquée à un jeu de données concret.
  \item \textbf{\faIcon{seedling}~Ce qu'il faut savoir avant} — les seuls
        prérequis de la section, et où les réviser.
  \item \textbf{\faIcon{exclamation-triangle}~Piège à éviter} — l'erreur
        classique que tout le monde fait.
  \item \textbf{\faIcon{key}~À retenir} — l'essentiel, si vous ne deviez
        garder qu'une chose.
  \item \textbf{\faIcon{rocket}~Pour aller plus loin} — un pont vers des
        sujets plus avancés, que vous pouvez ignorer en première lecture.
\end{itemize}

Les résultats importants sont énoncés comme \emph{théorèmes}, et le code
qui les met en œuvre (surtout en Python) est fourni et commenté. Les
démonstrations sont là pour ceux qui veulent savoir \emph{pourquoi} :
sautez-les sans remords à la première lecture, le fil du texte reste
compréhensible sans elles.

\vspace{1em}
\hfill\textit{Bonne lecture.}
\clearpage


\cleardoublepage
\pdfbookmark[0]{Table des matières}{toc}
\tableofcontents

%----------------------------------------------------------------------
% Corps du livre
%----------------------------------------------------------------------
\mainmatter

% Astuce : \input insère le chapitre directement. Pour de très gros livres,
% vous pouvez remplacer \input par \include (compilation chapitre par
% chapitre via \includeonly), à condition de n'utiliser QUE \include.
%
% Progression voulue : on ne suppose rien au départ, et chaque partie ne
% s'appuie que sur les précédentes. Décommentez les chapitres au fur et à
% mesure de leur rédaction.

\part{Le langage des données : algèbre linéaire}
\chapter{Vecteurs, matrices et normes}
\label{chap:vecteurs}

\begin{quote}\itshape
Avant de faire apprendre quoi que ce soit à une machine, apprivoisons les
objets avec lesquels elle travaille : les vecteurs et les matrices. Ce
chapitre pose le décor — en douceur.
\end{quote}

\begin{prerequis}
Savoir additionner et multiplier des nombres, et savoir ce qu'est une
racine carrée. Rien d'autre. Si le symbole $\sum$ vous inquiète, ne vous
arrêtez pas : il est expliqué la première fois qu'il apparaît.
\end{prerequis}

\section{Un vecteur, c'est juste une liste de nombres}

\begin{enclair}
Un \emph{vecteur} de taille $n$, c'est une liste ordonnée de $n$ nombres.
Rien de plus. Par exemple, la position d'un point sur une carte
$(x, y)$ est un vecteur de taille~2. Les pixels d'une image, mis bout à
bout, forment un vecteur (très long).
\end{enclair}

\begin{enia}
En intelligence artificielle, \textbf{tout} devient un vecteur : une
photo, une phrase, un client, une molécule. C'est la première chose que
fait n'importe quel système d'IA — transformer le monde en listes de
nombres, parce que c'est la seule chose qu'une machine sache calculer.
Le vecteur qui représente un objet s'appelle son \emph{plongement}
(\emph{embedding} en anglais).
\end{enia}

\begin{graphe}{Un jeu de données, c'est un nuage de vecteurs. Ici, chaque
fleur est décrite par deux nombres seulement — donc par un vecteur de
$\R^2$ —, et l'espèce se lit presque à l'œil nu.}[nuage-iris]
\begin{axis}[mbaxis, height=5.6cm,
             xlabel={Longueur du pétale (cm)},
             ylabel={Largeur du pétale (cm)},
             xmin=0, xmax=7.5, ymin=0, ymax=3,
             legend pos=north west]
  \addplot[only marks, mark=*,       mark size=1.5pt, mbBlue] coordinates {
    (1.4,0.2) (1.4,0.2) (1.3,0.2) (1.5,0.2) (1.4,0.2) (1.7,0.4)
    (1.4,0.3) (1.5,0.2) (1.4,0.2) (1.5,0.1) (1.5,0.2) (1.6,0.2)};
  \addplot[only marks, mark=triangle*, mark size=2pt,  mbAmber] coordinates {
    (4.7,1.4) (4.5,1.5) (4.9,1.5) (4.0,1.3) (4.6,1.5) (4.5,1.3)
    (4.7,1.6) (3.3,1.0) (4.6,1.3) (3.9,1.4) (4.2,1.5) (4.0,1.0)};
  \addplot[only marks, mark=square*, mark size=1.5pt, mbTeal] coordinates {
    (6.0,2.5) (5.1,1.9) (5.9,2.1) (5.6,1.8) (5.8,2.2) (6.6,2.1)
    (5.5,1.8) (6.3,1.8) (5.8,1.6) (6.1,2.5) (5.1,2.0) (5.3,1.9)};
  \legend{Espèce A, Espèce B, Espèce C}
\end{axis}
\end{graphe}

On note un vecteur en gras, $\vect{x} \in \Rn$, et ses composantes
$x_1, x_2, \dots, x_n$. On les empile verticalement :
\[
  \vect{x} = \begin{pmatrix} x_1 \\ x_2 \\ \vdots \\ x_n \end{pmatrix}.
\]

\begin{intuition}
Voyez le vecteur $\vect{x}\in\R^2$ comme une \emph{flèche} partant de
l'origine et pointant vers le point $(x_1, x_2)$. Additionner deux vecteurs,
c'est mettre les flèches bout à bout.
\end{intuition}

\begin{schema}{Le vecteur $\vect{x} = (3,4)$ vu comme une flèche. Sa
longueur se lit directement avec le théorème de Pythagore : les pointillés
gris forment un triangle rectangle de côtés $3$ et $4$, dont l'hypoténuse
— la flèche — mesure $5$.}[vecteur-fleche]
  \begin{scope}[x=1cm, y=1cm]
    \axesxy{5.2}{5.2}
    \draw[construction] (3,0) -- (3,4);
    \draw[construction] (0,4) -- (3,4);
    \draw[vecteur] (0,0) -- (3,4);
    \node[point] at (3,4) {};
    \node[etiquette, below] at (3,-0.12) {$x_1 = 3$};
    \node[etiquette, left]  at (-0.12,4) {$x_2 = 4$};
    \node[font=\small, above left=-2pt] at (1.5,2) {$\vect{x}$};
    \node[annotation bulle] (mblong) at (4.3,1.6)
         {longueur\\ $\norm{\vect{x}}_2 = 5$};
    \draw[trait annotation] (mblong) to[bend right=18] (2.1,2.5);
  \end{scope}
\end{schema}

\begin{definition}[Norme euclidienne][norme-eucl]
La \emph{longueur} d'un vecteur $\vect{x}\in\Rn$, appelée \emph{norme
euclidienne}, est le nombre
\[
  \norm{\vect{x}}_2 = \sqrt{x_1^2 + x_2^2 + \dots + x_n^2}
                    = \sqrt{\smash{\textstyle\sum_{i=1}^{n} x_i^2}}.
\]
C'est exactement le théorème de Pythagore, généralisé à $n$ dimensions.
\end{definition}

\begin{exemple}[Calcul d'une norme]
Pour $\vect{x} = (3, 4)$, on a $\norm{\vect{x}}_2 = \sqrt{9 + 16}
= \sqrt{25} = 5$. La flèche mesure 5~unités de long.
\end{exemple}

\begin{piege}
La norme n'est \textbf{pas} la somme des composantes ! On élève au carré,
on additionne, \emph{puis} on prend la racine. Oublier la racine est
l'erreur numéro~un des débutants.
\end{piege}

\section{Une matrice transforme les vecteurs}

Une matrice $A \in \Mmn{m}{n}$ est un tableau de nombres à $m$ lignes et
$n$ colonnes. Son vrai rôle : \textbf{transformer} un vecteur en un autre,
par le produit $\vect{y} = A\vect{x}$.

\begin{aretenir}
Multiplier par une matrice, c'est appliquer une transformation
géométrique : une rotation, un étirement, une projection\dots{} Comprendre
\emph{quelle} transformation cache une matrice, c'est comprendre l'algèbre
linéaire.
\end{aretenir}

\begin{enia}
Une couche de réseau de neurones n'est rien d'autre que cela : une
multiplication par une matrice, $\vect{y} = W\vect{x} + \vect{b}$, suivie
d'une petite fonction appliquée nombre par nombre. Les coefficients de $W$
sont précisément ce que le réseau \emph{apprend}. Entraîner un modèle,
c'est chercher les bons nombres à mettre dans des matrices — et rien de
plus mystérieux que cela.
\end{enia}

\begin{schema}{Un petit réseau de neurones. Chaque trait porte un
coefficient : ce sont exactement les cases de la matrice $W$. Un réseau
qui « apprend », c'est un réseau qui ajuste ces nombres.}[reseau-couche]
  \reseauneurones[Entrée $\vect{x}$][Couche cachée][Sortie $\vect{y}$]{3}{4}{2}
\end{schema}

\begin{theoreme}[Existence et unicité de la solution][inversible]
Soit $A \in \Mn$ une matrice carrée. Le système linéaire $A\vect{x} =
\vect{b}$ admet une \emph{unique} solution pour tout $\vect{b}\in\Rn$ si et
seulement si $A$ est \emph{inversible}, c'est-à-dire $\det(A) \neq 0$.
Dans ce cas, $\vect{x} = \inv{A}\,\vect{b}$.
\end{theoreme}

\begin{proof}
Si $A$ est inversible, poser $\vect{x} = \inv{A}\vect{b}$ donne
$A\vect{x} = A\inv{A}\vect{b} = \identity\vect{b} = \vect{b}$, d'où
l'existence ; et si $A\vect{x} = A\vect{x}'$ alors, en multipliant à gauche
par $\inv{A}$, on obtient $\vect{x} = \vect{x}'$, d'où l'unicité.
Réciproquement, si $A$ n'est pas inversible, son noyau $\Ker A$ contient un
vecteur non nul et la solution, quand elle existe, n'est jamais unique.
\end{proof}

\begin{remarque}
En pratique, on ne calcule \emph{presque jamais} $\inv{A}$ pour résoudre
$A\vect{x}=\vect{b}$ : c'est coûteux et numériquement instable. On préfère
des méthodes directes (l'élimination de Gauss, au chapitre suivant) ou
itératives. Voir aussi~\textcite{golubvanloan2013}.
\end{remarque}

\section{Passons au code}

Voici comment manipuler tout cela en Python avec \co{python}{numpy}, la
bibliothèque de référence en calcul scientifique~\parencite{harris2020numpy}.

\begin{code}{python}[Norme et résolution d'un système linéaire]
import numpy as np

# Un vecteur et sa norme euclidienne
x = np.array([3.0, 4.0])
print(np.linalg.norm(x))          # -> 5.0

# Résoudre A x = b SANS inverser A (recommandé)
A = np.array([[2.0, 1.0],
              [1.0, 3.0]])
b = np.array([1.0, 2.0])

x = np.linalg.solve(A, b)         # méthode directe, stable
print("solution :", x)
print("residu   :", np.linalg.norm(A @ x - b))
\end{code}

Le même calcul en Julia, très prisé en recherche numérique :

\begin{code}{julia}[La même résolution en Julia]
using LinearAlgebra

A = [2.0 1.0; 1.0 3.0]
b = [1.0, 2.0]

x = A \ b            # l'opérateur « backslash » résout le système
println("solution : ", x)
println("norme    : ", norm(x))
\end{code}

\begin{plusloin}
Le nombre $\cond(A) = \norm{A}\,\norm{\inv{A}}$, appelé
\emph{conditionnement}, mesure à quel point la solution est sensible aux
erreurs d'arrondi. Un grand conditionnement annonce des ennuis numériques.
Nous y reviendrons en détail.
\end{plusloin}

\section*{Exercices}
\addcontentsline{toc}{section}{Exercices}

\begin{exercice}[Norme et distance]
Calculez $\norm{\vect{x}}_2$ pour $\vect{x} = (1, 2, 2)$, puis la distance
entre $\vect{x}$ et $\vect{y} = (1, 0, 0)$.
\end{exercice}

\begin{solution}
$\norm{\vect{x}}_2 = \sqrt{1 + 4 + 4} = 3$. La distance vaut
$\norm{\vect{x}-\vect{y}}_2 = \norm{(0,2,2)}_2 = \sqrt{8} = 2\sqrt{2}$.
\end{solution}

\begin{exercice}[Inversibilité]
La matrice $A = \begin{psmallmatrix} 1 & 2 \\ 2 & 4 \end{psmallmatrix}$
est-elle inversible ? Justifiez, puis dites ce que cela implique pour le
système $A\vect{x} = \vect{b}$.
\end{exercice}

\begin{solution}
$\det(A) = 1\cdot 4 - 2\cdot 2 = 0$ : la matrice \textbf{n'est pas}
inversible (la 2\ieme{} ligne est le double de la 1\iere{}). D'après le
\cref{thm:inversible}, le système n'a pas de solution unique : soit aucune
solution, soit une infinité.
\end{solution}

% \input{chapters/02-systemes-lineaires}       % résoudre Ax = b, Gauss
% \input{chapters/03-factorisations}           % LU, Cholesky, QR
% \input{chapters/04-valeurs-propres-et-svd}   % ACP, compression, embeddings

\part{Le langage de l'incertain : probabilités et statistique}
% \input{chapters/05-probabilites}             % hasard, variables aléatoires
% \input{chapters/06-lois-usuelles}            % Bernoulli, normale, ...
% \input{chapters/07-estimation}               % moyenne, variance, maximum de vraisemblance
% \input{chapters/08-bayes}                    % probabilités conditionnelles, Bayes

\part{Bouger dans la bonne direction : calcul différentiel et optimisation}
% \input{chapters/09-derivees-et-gradients}
% \input{chapters/10-descente-de-gradient}
% \input{chapters/11-convexite}
% \input{chapters/12-optimisation-sous-contraintes}

\part{Apprendre à partir de données : apprentissage automatique}
% \input{chapters/13-regression-lineaire}
% \input{chapters/14-classification}           % régression logistique
% \input{chapters/15-generalisation}           % sur-apprentissage, régularisation

\part{Les réseaux de neurones, démontés pièce par pièce}
% \input{chapters/16-le-neurone}
% \input{chapters/17-reseaux-multicouches}
% \input{chapters/18-retropropagation}
% \input{chapters/19-entrainement-en-pratique}

%----------------------------------------------------------------------
% Pages finales
%----------------------------------------------------------------------
\backmatter
\cleardoublepage
\printbibliography[heading=bibintoc]

\end{document}
